\scanpage{315}
\mspage{14}
dans $\widehat{\boldsymbol\Delta}_{ab}=\mathrm{Ab}_*$ est \uline{universel} pour les
complexes de cochaînes dans des cat.
additives $A$ avec $\varinjlim$. \unsure{\&} (\uline{NB} On \unsure{a}
on peut aussi interpréter $\mathrm{Ab}_*$ comme $\mathrm{Ab}_\cdot$
dans cette description de la cat. contenant
l'objet universel --- ).

b) En dualisant et interprétant
$n\unsure{\vdash|\mapsto}\mathbb C_*({\Delta^n}_*)$ comme objet \uline{simplicial}
dans $(\mathrm{Ab}_*)^o\inserted{=(\mathrm{Ab}^o)^*}$, qui est une cat. additive
avec $\varprojlim$, il est universel pour les
cat. $B$ avec $\varprojlim$ et foncteurs entre
icelles y commutant. De même pour le
complexe de chaînes associé.

Du point de vue de l'\unsure{Algèbre} universelle,
\deleted{\unsure{cela veut dire que les constructions}}
\inserted{on peut exprimer le contenu de ces}
réflexions de la façon suivante.

1. Les inv. contravariants ``commutant aux
$\varprojlim$'' en $X_*\in\mathrm{Ss}$, à valeurs dans une
cat. additive $A$ avec $\varprojlim$, s'expriment
ou encore en termes d'objets simpliciaux
de $A$, ou de complexes de chaînes de $A$.

2. Les constructions qu'on peut faire sur des
tels invariants (\inserted{pour $A$ variable,} de façon compatible aux
foncteurs $A\longrightarrow A'$ commutant aux $\varprojlim$)
s'expriment à l'aide des objets
de $(\mathrm{Ab}_*)^o$ (où se trouve l'objet
simplicial --- pour cat. add. à $\varprojlim$ ---
\uline{universel}), ou de $(\mathrm{Ab}_\cdot)$

et $1')$, $2')$ les énoncés duaux (formulés
à partir de $1)$, $2)$)

Si $F,K_*,K_\cdot$ se correspondent dans
1, on a
\[
\left\{\begin{aligned}
F(X_*)&\simeq\operatorname{Hom}_*(\mathbb C_*(X_*),K_*)\simeq
\operatorname{Hom}_\cdot(\mathbb C_\cdot(X_*),K_\cdot)\\
K_\cdot&=\mathrm{ND}_\cdot(K_*),\qquad K_*=\mathrm{DP}_*(K_\cdot)
\end{aligned}\right.
\]
Si de même $\Omega$ est une ``opération sur des
invariants'', et si $L_*\inserted{\in\mathrm{Ab}_*}$ et $L_\cdot\in\mathrm{Ab}_\cdot$ y
sont associés, on aura
\[
\left\{\begin{aligned}
\Omega(F)&\simeq\operatorname{Hom}_*(L_*,K_*)\overset{\inserted{=\Omega_*(K_*)}}\simeq
\operatorname{Hom}_\cdot(L_\cdot,K_\cdot)\inserted{=\Omega_\cdot(K_\cdot)}\\
L_\cdot&=\mathrm{ND}_*(L_*),\qquad L_*=\mathrm{DP}_*(L_\cdot)
\end{aligned}\right.
\]
\margin{$\Omega$ a deux interprétations ``opérations sur des $K^*$'' ($\Omega^*$) et ``op. sur des $K^\cdot$'' ($\Omega^\cdot$) --- pour trouver des objets de $B$ \deleted{\unsure{provenant}} à partir d'un complexe (simp. ou de chaînes) dans $B$}
(La signification des symboles $\operatorname{Hom}_*$ (resp $\operatorname{Hom}_\cdot$)
pour tt un $C_*$ ou $L_*$ \inserted{dans $\mathrm{Ab}_*$} (resp. $C_\cdot$ ou $L_\cdot$ \inserted{dans $\mathrm{Ab}_\cdot$})
a été expliquée plus haut).

Dualement, \inserted{pour} les \deleted{\unsure{opérations}}\inserted{invariants} \inserted{d'objets $X_*\in\mathrm{Ss}$} covariants
commutant aux $\varinjlim$, à valeurs dans
$A$, si $G,K^*,K^\cdot$ se correspondent,
on aura
\[
\left\{\begin{aligned}
G(X_*)&\simeq\deleted{\unsure{Z^0}}(\mathbb C_*(X_*)\mathbin{\overset{*}{\otimes}}\unsure{{K_*}^*|K^*})
\simeq\deleted{\unsure{Z^0}}(\mathbb C_\cdot(X_*)\mathbin{\overset{\cdot}{\otimes}}K^\cdot)\\
\unsure{K_\cdot|K^\cdot}&=\mathrm{ND}^{\cdot}(K^*),\qquad K^*=\mathrm{ND}^*(K^\cdot)
\end{aligned}\right.
\]
\margin{à préciser la signification de ces formules}
et si l'opération $\Omega$, $L_*$ et $L_\cdot$ se
correspondent, ($L_*\in\mathrm{Ab}_*$, $L_\cdot\in\mathrm{Ab}_\cdot$)
on aura
\[
\left\{\begin{aligned}
\Omega(G)&\simeq L_*\mathbin{\overset{*}{\otimes}}K^*\simeq L_\cdot\mathbin{\overset{\cdot}{\otimes}}K^\cdot\\
L_\cdot&\simeq\mathrm{ND}_\cdot(L_*),\qquad L_*\simeq\mathrm{DP}_*(L_\cdot)
\end{aligned}\right.
\]
\margin{316}
